%%%%%%%%%%%%
%
% $Autor: Wings $
% $Datum: 2019-03-05 08:03:15Z $
% $Pfad: Functions.tex $
% $Version: 4250 $
% !TeX spellcheck = en_GB/de_DE
% !TeX encoding = utf8
% !TeX root = manual 
% !TeX TXS-program:bibliography = txs:///biber
%
%%%%%%%%%%%%

\chapter{Bedienung und Anzeige}

\section{Das Radar-Bild interpretieren}
Auf Ihrem Bildschirm sehen Sie eine kreisförmige Anzeige.
\begin{itemize}
	\item \textbf{Grüne Linie:} Zeigt die aktuelle Blickrichtung des Radars.
	\item \textbf{Rote Markierungen:} Hier wurde ein Hindernis erkannt. Je näher der Punkt an der Mitte ist, desto näher ist der Gegenstand.
	\item \textbf{Zahlenwerte:} Am Rand wird Ihnen die Entfernung in Zentimetern (cm) angezeigt.
\end{itemize}

\section{Reichweite}
Das System erkennt Objekte zuverlässig in einer Distanz von ca. $2\,cm$ bis $400\,cm$. Sehr weiche Oberflächen (wie Vorhänge) können den Schall schlucken und werden evtl. später erkannt als harte Oberflächen (wie Wände).
